\chapter{从内积到 Hilbert 空间}
\label{chap:math-hilbert}

有限维线性代数里，正交投影总能找到，Cauchy 列也总有极限。波函数却是函数而非有限列数字。若用函数列逼近一个态，我们需要先回答两件事：怎样判断它们越来越近，以及极限是否仍是允许的态。

\section{函数空间中“越来越近”是什么意思}

\begin{definition}[度量、收敛与完备]
集合 \(X\) 上的度量 \(d:X\times X\to[0,\infty)\) 满足 \(d(x,y)=0\iff x=y\)、\(d(x,y)=d(y,x)\) 和三角不等式。若对每个 \(\varepsilon>0\) 存在 \(N\)，使 \(n\ge N\) 时 \(d(x_n,x)<\varepsilon\)，称 \(x_n\to x\)。若对每个 \(\varepsilon>0\) 存在 \(N\)，使 \(m,n\ge N\) 时 \(d(x_m,x_n)<\varepsilon\)，称 \((x_n)\) 为 Cauchy 列。每个 Cauchy 列都在 \(X\) 内收敛时，\(X\) 完备。
\end{definition}

在有理数中取逐位截断 \(\sqrt2\) 的小数列，它是 Cauchy 列，却没有有理数极限。这说明“项彼此接近”不足以保证极限留在选定的集合。对函数还要指定距离：\(f_n(x)=x^n\) 在 \([0,1]\) 的积分平方距离下趋于零，但 \(\sup_{[0,1]}|f_n|=1\)，因而不一致趋于零。同一个函数列在不同距离下可以有不同结论。

\begin{definition}[范数、内积与 Hilbert 空间]
复向量空间上的范数 \(\|\cdot\|\) 满足正定、\(\|\lambda x\|=|\lambda|\|x\|\) 和三角不等式。内积 \(\langle x,y\rangle\) 对第二变量线性，满足 \(\langle x,y\rangle=\overline{\langle y,x\rangle}\) 与 \(\langle x,x\rangle>0\)（\(x\ne0\)）。它诱导范数 \(\|x\|=\sqrt{\langle x,x\rangle}\)。对此范数完备的内积空间称 Hilbert 空间。
\end{definition}

内积确实给出范数所需的三角不等式，关键是下一个估计。

\begin{theorem}[Cauchy--Schwarz 不等式]
在任意实或复内积空间中，\(|\langle x,y\rangle|\le\|x\|\|y\|\)。等号成立当且仅当 \(x,y\) 线性相关。
\end{theorem}
\begin{proof}
若 \(y=0\)，结论成立。否则取 \(\lambda=\langle y,x\rangle/\|y\|^2\)。由范数非负及内积第二变量线性，
\[
0\le\|x-\lambda y\|^2
=\|x\|^2-\lambda\langle x,y\rangle
-\overline\lambda\langle y,x\rangle+|\lambda|^2\|y\|^2
=\|x\|^2-\frac{|\langle x,y\rangle|^2}{\|y\|^2}.
\]
移项得到不等式。等号等价于 \(x-\lambda y=0\)。
\end{proof}

这里先不预设测度论，也能严格得到后文需要的 \(L^2(0,1)\)。
在连续函数空间 \(C[0,1]\) 上取
\(\|f\|_2^2=\int_0^1|f(x)|^2\,\dd x\)。
把它的所有 Cauchy 列放在一起，若两列 \((f_n),(g_n)\) 满足
\(\|f_n-g_n\|_2\to0\)，就把它们视为同一个元素；
所得完备化定义为 \(L^2(0,1)\)。加法、数乘与内积均取
连续函数列的极限，刚证出的 Cauchy--Schwarz 不等式保证内积极限
存在且不依赖代表列。为验证完备性，给定完备化中的 Cauchy 列，
先抽子列使相邻两项距离小于 \(2^{-j}\)，再从第 \(j\) 项
选连续函数代表使误差小于 \(2^{-j}\)。三角不等式说明这些
连续函数构成 Cauchy 列；它所代表的元素是抽出子列的极限，
也就是原 Cauchy 列的极限。

在通常的函数实现中，\(L^2\) 的元素可看作平方可积函数，
但两个函数只要在零测集以外相等，就代表同一个元素。
这里“零测集”指对任意 \(\varepsilon>0\)，都能用可数个开区间
覆盖它，且这些区间总长度小于 \(\varepsilon\) 的集合；
“几乎处处相等”就是除这样的集合外处处相等。
例如只在一个点上改动函数值不改变其 \(L^2\) 元素。
其内积写成 \(\langle f,g\rangle=
\int_0^1\overline{f(x)}g(x)\,\dd x\)，这里的零向量
是一整个几乎处处为零的函数等价类。只取多项式并沿用这个范数
则不完备：多项式可在平方平均意义下逼近阶跃函数，而极限不是多项式。

\section{正交投影的存在与唯一性}

有限维里解最小二乘方程时，我们把向量投到一列向量张成的子空间。无穷维中，距离的下确界未必由子空间里的元素取得；这里第一次真正用到完备性和闭性。

\begin{definition}[闭子空间与正交补]
子空间 \(M\subset\mathcal H\) 称为闭的，若 \(M\) 中任意收敛列的极限仍属于 \(M\)。它的正交补为 \(M^\perp=\{z\in\mathcal H:\langle m,z\rangle=0\text{ 对每个 }m\in M\}\)。
\end{definition}

例如 \(L^2(0,1)\) 中的零均值函数构成闭子空间，因为积分对 \(L^2\) 范数连续。相反，多项式集合不是闭的；点到该集合的距离可能为零，但点不属于集合，这时谈不上最近的多项式。

\begin{theorem}[闭子空间投影定理]
若 \(M\) 是 Hilbert 空间 \(\mathcal H\) 的闭线性子空间，则每个 \(x\in\mathcal H\) 恰有一个 \(p\in M\) 使 \(\|x-p\|=\inf_{m\in M}\|x-m\|\)。该元素满足 \(x-p\in M^\perp\)，故 \(x=p+(x-p)\) 是唯一的正交分解。
\end{theorem}
\begin{proof}
令距离下确界为 \(d\)，选 \(m_n\in M\) 使 \(\|x-m_n\|^2\to d^2\)。平行四边形恒等式用于 \(x-m_n\) 和 \(x-m_k\)，给出
\[
\|m_n-m_k\|^2
=2\|x-m_n\|^2+2\|x-m_k\|^2
-4\left\|x-\frac{m_n+m_k}{2}\right\|^2
\le 2\|x-m_n\|^2+2\|x-m_k\|^2-4d^2.
\]
右端趋于零；这里用到了中点仍在 \(M\)。故 \((m_n)\) 为 Cauchy 列，由 \(\mathcal H\) 完备且 \(M\) 闭，存在 \(p\in M\) 使 \(m_n\to p\)。范数连续，所以 \(\|x-p\|=d\)。

现在任取 \(h\in M\)，实数 \(t\) 的二次式 \(\|x-p-th\|^2\) 在 \(t=0\) 取最小值，故 \(\operatorname{Re}\langle x-p,h\rangle=0\)。将 \(h\) 换成 \(\ii h\) 又得到虚部为零，因此 \(x-p\perp M\)。若 \(q\in M\) 也是最小点，则 \(x-q\perp M\)，从而 \(p-q\in M\cap M^\perp=\{0\}\)。
\end{proof}

记 \(P_Mx=p\)。由正交分解，\(\|x-m\|^2=\|x-P_Mx\|^2+\|P_Mx-m\|^2\)；这个式子既证明最优性，也说明投影把误差分成不能消除的部分与选择 \(m\) 额外造成的部分。对常数函数子空间，\(P_Mf=\int_0^1f(x)\,\dd x\)，因为 \(f-P_Mf\) 与每个常数正交。

\section{连续线性测量可以写成内积}

对波函数取加权积分 \(F(f)=\int_0^1\overline{g(x)}f(x)\,\dd x\)，得到一个复数。它对 \(f\) 线性，且由 Cauchy--Schwarz 不等式有 \(|F(f)|\le\|g\|_2\|f\|_2\)。这类测量称为连续线性泛函；连续性意味着输入的小扰动不会造成有限大小的输出跳变。

\begin{definition}[线性泛函及其范数]
\(F:\mathcal H\to\mathbb C\) 对加法和复数乘法保持线性，称线性泛函。若存在 \(C<\infty\) 使 \(|F(x)|\le C\|x\|\) 对所有 \(x\) 成立，称它有界，并定义 \(\|F\|=\sup_{\|x\|\le1}|F(x)|\)。
\end{definition}

\begin{theorem}[Riesz 表示定理]
在复 Hilbert 空间 \(\mathcal H\) 上，每个有界线性泛函 \(F\) 都唯一地写成 \(F(x)=\langle y,x\rangle\)，其中 \(y\in\mathcal H\)，并且 \(\|F\|=\|y\|\)。
\end{theorem}
\begin{proof}
若 \(F=0\)，取 \(y=0\)。否则 \(N=\ker F\) 是闭子空间：若 \(x_n\in N\) 且 \(x_n\to x\)，则 \(|F(x)|\le\|F\|\|x-x_n\|\to0\)。由投影定理，取 \(v\in N^\perp\) 非零。对任意 \(x\)，\(x-F(x)v/F(v)\in N\)。又 \(N^\perp\) 只有一维：若 \(w\in N^\perp\)，同样的差既属于 \(N\) 又属于 \(N^\perp\)，故为零。于是
\[
F(x)=\frac{F(v)}{\|v\|^2}\langle v,x\rangle
=\left\langle\frac{\overline{F(v)}}{\|v\|^2}v,x\right\rangle.
\]
取括号中的向量为 \(y\)。若两个向量表示同一泛函，取 \(x\) 为其差即知二者相等。Cauchy--Schwarz 给 \(\|F\|\le\|y\|\)；当 \(y\ne0\) 取 \(x=y/\|y\|\) 得反向不等式。
\end{proof}

这个定理并不是说每个线性规则都由向量表示。例如在 \(L^2(0,1)\) 中“取 \(f(0)\)”对函数的几乎处处等价类根本没有定义；在连续函数上它虽有定义，却对 \(L^2\) 范数不连续。一个尖峰可以保持 \(f(0)=1\)，而平方积分趋于零。连续性是定理不可删的条件。

Riesz 定理告诉我们如何表示单个测量。若有一列两两正交、范数均为一的向量 \(e_1,e_2,\ldots\)，就可以同时记录许多测量 \(\langle e_n,x\rangle\)。但无限个系数是否足以重建 \(x\)，不能照搬有限维的“维数相同”。先对前 \(N\) 个方向作正交投影 \(P_Nx=\sum_{n=1}^N e_n\langle e_n,x\rangle\)。逐项取内积得到 \(x-P_Nx\perp P_Nx\)，所以
\[
 \|x\|^2=\sum_{n=1}^N|\langle e_n,x\rangle|^2+\|x-P_Nx\|^2.
\]
右端最后一项非负，令 \(N\) 增大便得到 Bessel 不等式 \(\sum_n|\langle e_n,x\rangle|^2\le\|x\|^2\)。特别地，系数的平方和有限，且对 \(M>N\)，\(\|P_Mx-P_Nx\|^2=\sum_{n=N+1}^M|\langle e_n,x\rangle|^2\to0\)。Hilbert 空间完备性使 \(P_Nx\) 收敛到某个 \(p\)。剩余向量 \(x-p\) 与每个 \(e_n\) 正交；它不必为零，除非这些向量的有限线性组合稠密。这后一条件才叫作\emph{完备正交系}。

\begin{theorem}[正交展开与 Parseval 恒等式]
 对一组标准正交向量 \((e_n)\)，以下三件事等价：其有限线性组合在 \(\mathcal H\) 中稠密；与每个 \(e_n\) 正交的向量只能为零；每个 \(x\in\mathcal H\) 都满足
 \begin{equation}
 x=\sum_{n=1}^\infty e_n\langle e_n,x\rangle,
 \qquad \|x\|^2=\sum_{n=1}^\infty|\langle e_n,x\rangle|^2.
 \label{eq:math-parseval}
 \end{equation}
 级数在 Hilbert 空间范数下收敛。
\end{theorem}
\begin{proof}
 上一段已证明 \(P_Nx\to p\)，且 \(x-p\) 与每个 \(e_n\) 正交。若有限线性组合稠密，它也与全空间正交，故为零；若第二个条件成立，也直接得到 \(x=p\)。取有限维勾股等式的极限即得 Parseval。反过来，若展开式对每个 \(x\) 成立，每个向量都是有限线性组合的范数极限，故这些组合稠密；若只知第二个条件成立，上述 \(x-p\) 的论证已经证明展开式。故三者等价。
\end{proof}

对后文的 Dirichlet 算子，需要的不只是任意完备系，而是正弦函数确实完备。这里给出一个不依赖“傅里叶级数显然收敛”的证明。取 \(g\in C_c(0,1)\)，先把它奇延拓到 \((-1,1)\)，再作周期二的连续延拓 \(G\)。对周期二连续函数，Fejér 核
\[
 F_N(t)=\frac1{N+1}\left(\frac{\sin((N+1)\pi t/2)}{\sin(\pi t/2)}\right)^2
\]
非负且在一个周期上的积分为 \(2\)。它也等于 \((N+1)^{-1}|\sum_{j=0}^N\ee^{\ii\pi jt}|^2\)，所以积分值由有限项正交性直接算出。若 \(0<\delta<1\)，在 \(\delta\le|t|\le1\) 有 \(F_N(t)\le[(N+1)\sin^2(\pi\delta/2)]^{-1}\)。因此卷积平均 \(\frac12\int_{-1}^1F_N(t)G(x-t)\,\dd t\) 与 \(G(x)\) 的差，可以拆成 \(|t|<\delta\) 上由一致连续性控制的部分，以及余区间上随 \(N\to\infty\) 消失的部分；它一致趋于 \(G\)。该平均是有限三角多项式，且 \(G\) 为奇函数，所以只含正弦项。又 \(C_c(0,1)\) 在 \(L^2(0,1)\) 中稠密：先用连续函数作 \(L^2\) 逼近，再用靠近两端的截断函数乘它，截掉部分的平方积分随截断区间长度趋零。于是 \(\sqrt2\sin(n\pi x)\)、\(n\ge1\)，是 \(L^2(0,1)\) 的完备正交系。后文写 Fourier 正弦展开和 Parseval 时，所需的完备性已经在这里得到证明。

\section{有界算子与紧算子}

泛函输出一个数；算子可以输出另一个函数。对 Hilbert 空间间的线性算子 \(T:\mathcal H\to\mathcal K\)，若 \(\|Tx\|\le C\|x\|\) 对所有 \(x\) 成立，称有界，其最小可行 \(C\) 为 \(\|T\|\)。线性算子有界当且仅当它在零点连续：前向由估计直接得出；反向由连续性给出一个 \(\delta>0\)，使 \(\|x\|<\delta\) 时 \(\|Tx\|<1\)，再把任意非零 \(x\) 缩放到 \(\delta/2\) 即得 \(\|Tx\|\le2\|x\|/\delta\)。

有限维有界集的闭包紧，无穷维不再如此。为了恢复一部分矩阵的谱性质，需要更强条件。

\begin{definition}[紧算子]
有界线性算子 \(K:\mathcal H\to\mathcal K\) 若把每个有界序列映成含收敛子列的序列，称为紧算子。
\end{definition}

在 \(\ell^2\) 中，\(K(x_1,x_2,\ldots)=(x_1,x_2/2,\ldots)\) 是紧的。令 \(K_N\) 只保留前 \(N\) 项，则 \(K_N\) 的像在有限维空间中有收敛子列，且 \(\|K-K_N\|\le1/(N+1)\)。依次选子列并用这个一致误差，可得 \(K\) 紧。恒等算子却不紧，因为标准正交列 \((e_n)\) 中任意两项距离都是 \(\sqrt2\)，不存在 Cauchy 子列。

\begin{theorem}[紧自伴算子的谱结构]
若 \(K\) 是复 Hilbert 空间上的紧自伴算子，则它的非零谱点均为实本征值，各自具有有限重数，非零本征值只能向零累积；\(\mathcal H\) 有由这些本征向量和 \(\ker K\) 的正交基组成的正交基。
\end{theorem}

\begin{proof}
先证 \(K\ne0\) 时存在非零本征值。记 \(r=\|K\|^2>0\)，选单位向量 \(x_j\) 使 \(\|Kx_j\|^2\to r\)。令 \(L=K^2\)。由于 \(K\) 自伴，\(\langle x,Lx\rangle=\|Kx\|^2\)，且
\[
\|Lx\|^2=\|K(Kx)\|^2\le r\|Kx\|^2=r\langle x,Lx\rangle.
\]
因此
\[
\|(L-rI)x_j\|^2
=\|Lx_j\|^2-2r\langle x_j,Lx_j\rangle+r^2
\le r\bigl(r-\langle x_j,Lx_j\rangle\bigr)\longrightarrow0.
\]
紧算子 \(K\) 的平方 \(L\) 仍紧，所以 \((Lx_j)\) 有收敛子列；沿此子列，\(x_j=r^{-1}(Lx_j-(L-rI)x_j)\) 也收敛到单位向量 \(x\)，并有 \(K^2x=rx\)。于是 \(K(x\pm Kx/\sqrt r)=\pm\sqrt r\,(x\pm Kx/\sqrt r)\)。这两个向量至少一个非零，给出实本征值 \(\pm\|K\|\)。

取一个单位本征向量 \(e_1\)。因 \(K\) 自伴，\(e_1^\perp\) 在 \(K\) 下不变；限制在这个闭子空间上的算子仍紧且自伴。若它不为零，就用上段方法取 \(e_2\)；反复得到两两正交的 \(e_j\)，相应 \(\lambda_j\ne0\) 按绝对值不增排列。若过程无限延续而 \(|\lambda_j|\not\to0\)，则 \((Ke_j)=(\lambda_je_j)\) 没有收敛子列，与紧性矛盾。任一固定非零本征值也不能有无限维本征空间，理由相同。

令 \(M\) 为全部 \(e_j\) 的闭线性张成空间的正交补。它仍在 \(K\) 下不变。若 \(K|_M\ne0\)，首段会在 \(M\) 中给出绝对值为 \(a>0\) 的本征值；每次选取极值时的 \(|\lambda_j|\) 均不小于 \(a\)，与 \(\lambda_j\to0\) 矛盾。故 \(K|_M=0\)，也就是 \(M\subset\ker K\)。反向包含对非零本征向量的正交性亦成立，于是非零本征向量与 \(\ker K\) 的正交基张成整个 \(\mathcal H\)。若上述迭代有限终止，同样由终止条件得到结论。最后，若 \(z\ne0\) 且不等于任一 \(\lambda_j\)，则 \(|\lambda_j-z|\) 有共同正下界；在已得正交分解上逐项除以 \(\lambda_j-z\)，在核上除以 \(-z\)，便构成有界的 \((K-zI)^{-1}\)。所以没有遗漏的非零谱点。
\end{proof}

在对角例 \(Ke_n=e_n/n\) 中，本征值 \(1/n\) 趋于零；证明里的紧性正是排除其他有限累积点的关键。微分算子通常连有界性都不满足，下一章不能直接套用这个定理。

\begin{exercise}
取 \(f_n(x)=\max(1-nx,0)\) 于 \([0,1]\)。证明 \(f_n(0)=1\) 而 \(\|f_n\|_2\to0\)，由此说明取点值不是 \(L^2\) 上的连续线性泛函。
\end{exercise}


\begin{exercise}
设 \(M=\{f\in L^2(0,1):\int_0^1f=0\}\)。求 \(P_Mf\)，检验所得差向量与 \(M\) 正交；再说明若把 \(M\) 换成稠密但非闭的多项式子空间，投影定理哪一步失效。
\end{exercise}
